\documentclass[10pt,a4paper]{amsart}
\usepackage{amsmath,amssymb,amsthm,graphicx}
\usepackage{verbatim}
\newenvironment{reference}{\comment}{\endcomment}
%\newenvironment{reference}{}{}
\newenvironment{slogan}{\comment}{\endcomment}
\newenvironment{history}{\comment}{\endcomment}

% For commutative diagrams we use Xy-pic
\usepackage[all]{xy}

% We use 2cell for 2-commutative diagrams.
\xyoption{2cell}
% Theorem environments.
%
\theoremstyle{plain}
\newtheorem{theorem}[subsection]{Theorem}
\newtheorem{proposition}[subsection]{Proposition}
\newtheorem{lemma}[subsection]{Lemma}

\theoremstyle{definition}
\newtheorem{definition}[subsection]{Definition}
\newtheorem{example}[subsection]{Example}
\newtheorem{exercise}[subsection]{Exercise}
\newtheorem{situation}[subsection]{Situation}

\theoremstyle{remark}
\newtheorem{remark}[subsection]{Remark}
\newtheorem{remarks}[subsection]{Remarks}

\numberwithin{equation}{subsection}

% Macros
%
\def\lim{\mathop{\mathrm{lim}}\nolimits}
\def\colim{\mathop{\mathrm{colim}}\nolimits}
\def\Spec{\mathop{\mathrm{Spec}}}
\def\Hom{\mathop{\mathrm{Hom}}\nolimits}
\def\Ext{\mathop{\mathrm{Ext}}\nolimits}
\def\SheafHom{\mathop{\mathcal{H}\!\mathit{om}}\nolimits}
\def\SheafExt{\mathop{\mathcal{E}\!\mathit{xt}}\nolimits}
\def\Sch{\mathit{Sch}}
\def\Mor{\mathop{\mathrm{Mor}}\nolimits}
\def\Ob{\mathop{\mathrm{Ob}}\nolimits}
\def\Sh{\mathop{\mathit{Sh}}\nolimits}
\def\NL{\mathop{N\!L}\nolimits}
\def\CH{\mathop{\mathrm{CH}}\nolimits}
\def\proetale{{pro\text{-}\acute{e}tale}}
\def\etale{{\acute{e}tale}}
\def\QCoh{\mathit{QCoh}}
\def\Ker{\mathop{\mathrm{Ker}}}
\def\Im{\mathop{\mathrm{Im}}}
\def\Coker{\mathop{\mathrm{Coker}}}
\def\Coim{\mathop{\mathrm{Coim}}}

% Boxtimes
%
\DeclareMathSymbol{\boxtimes}{\mathbin}{AMSa}{"02}

%
% Macros for moduli stacks/spaces
%
\def\QCohstack{\mathcal{QC}\!\mathit{oh}}
\def\Cohstack{\mathcal{C}\!\mathit{oh}}
\def\Spacesstack{\mathcal{S}\!\mathit{paces}}
\def\Quotfunctor{\mathrm{Quot}}
\def\Hilbfunctor{\mathrm{Hilb}}
\def\Curvesstack{\mathcal{C}\!\mathit{urves}}
\def\Polarizedstack{\mathcal{P}\!\mathit{olarized}}
\def\Complexesstack{\mathcal{C}\!\mathit{omplexes}}
% \Pic is the operator that assigns to X its picard group, usage \Pic(X)
% \Picardstack_{X/B} denotes the Picard stack of X over B
% \Picardfunctor_{X/B} denotes the Picard functor of X over B
\def\Pic{\mathop{\mathrm{Pic}}\nolimits}
\def\Picardstack{\mathcal{P}\!\mathit{ic}}
\def\Picardfunctor{\mathrm{Pic}}
\def\Deformationcategory{\mathcal{D}\!\mathit{ef}}
\usepackage{cleveref}
\crefname{originalsection}{original section}{original sections}
\Crefname{originalsection}{Original section}{Original sections}
\makeatletter
\newlabel{section-cycle-of-closed}{{0AZA}{source}{}{sourcepointer-1}{}}
\newlabel{section-cycle-of-closed@cref}{{[originalsection][0][]0AZA}{[1][0][]source}{}{}{}}
\makeatother
\crefname{originalreference}{original reference}{original references}
\Crefname{originalreference}{Original reference}{Original references}
\makeatletter
\newlabel{varieties-lemma-dimension-locally-algebraic}{{0A21}{source}{}{sourcepointer-2}{}}
\newlabel{varieties-lemma-dimension-locally-algebraic@cref}{{[originalreference][0][]0A21}{[1][0][]source}{}{}{}}
\makeatother
\crefname{originallemma}{original lemma}{original lemmas}
\Crefname{originallemma}{Original lemma}{Original lemmas}
\makeatletter
\newlabel{lemma-transversal}{{0B1I}{source}{}{sourcepointer-3}{}}
\newlabel{lemma-transversal@cref}{{[originallemma][0][]0B1I}{[1][0][]source}{}{}{}}
\makeatother
\crefname{originalreference}{original reference}{original references}
\Crefname{originalreference}{Original reference}{Original references}
\makeatletter
\newlabel{algebra-lemma-regular-ring-CM}{{00NQ}{source}{}{sourcepointer-4}{}}
\newlabel{algebra-lemma-regular-ring-CM@cref}{{[originalreference][0][]00NQ}{[1][0][]source}{}{}{}}
\makeatother
\crefname{originalreference}{original reference}{original references}
\Crefname{originalreference}{Original reference}{Original references}
\makeatletter
\newlabel{algebra-lemma-regular-domain}{{00NP}{source}{}{sourcepointer-5}{}}
\newlabel{algebra-lemma-regular-domain@cref}{{[originalreference][0][]00NP}{[1][0][]source}{}{}{}}
\makeatother
\crefname{originalreference}{original reference}{original references}
\Crefname{originalreference}{Original reference}{Original references}
\makeatletter
\newlabel{algebra-lemma-one-equation}{{00KW}{source}{}{sourcepointer-6}{}}
\newlabel{algebra-lemma-one-equation@cref}{{[originalreference][0][]00KW}{[1][0][]source}{}{}{}}
\makeatother
\title{Multiplicity one on subscheme cycles: coefficient one, smoothness and non-containment}
\author{The Stacks Project Authors}
\date{stacks-project:intersection.tex:local-snapshot-20261012; GNU FDL 1.2 or later}
\begin{document}
\maketitle
\noindent\textit{Editorial note.} All delivered mathematical text is the author's own, copied verbatim from the pinned Stacks Project chapter file intersection.tex; no statement or proof step was rewritten, reordered or removed.\par\medskip
\noindent\textit{Editorial note.} Span recomputation by label pairing (LANE\_BRIEF 10.2): in this chapter the author's proof environment follows the statement environment as a standalone block; the spec covers the complete environment pair (statement span, proof span) as recorded in the pool candidate.\par\medskip
\noindent\textit{Editorial note.} Occupation scan: combined claim\_key + statement\_anchor + generic-line scan over all registration files (active and withdrawn) found 0 hits for this chapter (is-occupied.json, 2026-10-14); consistent with the duty manager's six-check dispatch.\par\medskip
\ifdefined\hypertarget\hypertarget{sourcepointer-1}{}\fi
\noindent\textit{Source reference.} Original section 0AZA: sources/stacks-chow/tags.txt, source lines 13818--13818. Cross-chapter standard prerequisite cited by the author inside the counted span; it is used as a pointer and its own proof is not part of this item.\par\medskip
\ifdefined\hypertarget\hypertarget{sourcepointer-2}{}\fi
\noindent\textit{Source reference.} Original reference 0A21: sources/stacks-chow/tags.txt, source lines 12689--12689. Cross-chapter standard prerequisite cited by the author inside the counted span; it is used as a pointer and its own proof is not part of this item.\par\medskip
\ifdefined\hypertarget\hypertarget{sourcepointer-3}{}\fi
\noindent\textit{Source reference.} Original lemma 0B1I: sources/stacks-chow/tags.txt, source lines 13896--13896. Cross-chapter standard prerequisite cited by the author inside the counted span; it is used as a pointer and its own proof is not part of this item.\par\medskip
\ifdefined\hypertarget\hypertarget{sourcepointer-4}{}\fi
\noindent\textit{Source reference.} Original reference 00NQ: sources/stacks-chow/tags.txt, source lines 868--868. Cross-chapter standard prerequisite cited by the author inside the counted span; it is used as a pointer and its own proof is not part of this item.\par\medskip
\ifdefined\hypertarget\hypertarget{sourcepointer-5}{}\fi
\noindent\textit{Source reference.} Original reference 00NP: sources/stacks-chow/tags.txt, source lines 867--867. Cross-chapter standard prerequisite cited by the author inside the counted span; it is used as a pointer and its own proof is not part of this item.\par\medskip
\ifdefined\hypertarget\hypertarget{sourcepointer-6}{}\fi
\noindent\textit{Source reference.} Original reference 00KW: sources/stacks-chow/tags.txt, source lines 766--766. Cross-chapter standard prerequisite cited by the author inside the counted span; it is used as a pointer and its own proof is not part of this item.\par\medskip
\begin{lemma}
\label{lemma-transversal-subschemes}
Let $X$ be a nonsingular variety. Let $r, s \geq 0$ and let
$Y, Z \subset X$ be closed subschemes with $\dim(Y) \leq r$ and
$\dim(Z) \leq s$. Assume $[Y]_r = \sum n_i[Y_i]$ and
$[Z]_s = \sum m_j[Z_j]$ intersect properly.
Let $T$ be an irreducible component of $Y_{i_0} \cap Z_{j_0}$
for some $i_0$ and $j_0$ and assume that the multiplicity
(in the sense of Section \ref{section-cycle-of-closed}) of $T$
in the closed subscheme $Y \cap Z$ is $1$.
Then
\begin{enumerate}
\item the coefficient of $T$ in $[Y]_r \cdot [Z]_s$ is $1$,
\item $Y$ and $Z$ are nonsingular at the generic point of $T$,
\item $n_{i_0} = 1$, $m_{j_0} = 1$, and
\item $T$ is not contained in $Y_i$ or $Z_j$ for $i \not = i_0$ and
$j \not = j_0$.
\end{enumerate}
\end{lemma}

\begin{proof}
Set $n = \dim(X)$, $a = n - r$, $b = n - s$. Observe that
$\dim(T) = r + s - n = n - a - b$ by the assumption that the
intersections are proper. Let $(A, \mathfrak m, \kappa) =
(\mathcal{O}_{X, \xi}, \mathfrak m_\xi, \kappa(\xi))$ where $\xi \in T$
is the generic point. Then $\dim(A) = a + b$, see
Varieties, Lemma \ref{varieties-lemma-dimension-locally-algebraic}.
Let $I_0, I, J_0, J \subset A$ cut out the trace of
$Y_{i_0}$, $Y$, $Z_{j_0}$, $Z$ in $\Spec(A)$.
Then $\dim(A/I) = \dim(A/I_0) = b$ and $\dim(A/J) = \dim(A/J_0) = a$
by the same reference. Set $\overline{I} = I + \mathfrak m^2/\mathfrak m^2$.
Then $I \subset I_0 \subset \mathfrak m$ and
$J \subset J_0 \subset \mathfrak m$ and $I + J = \mathfrak m$.
By Lemma \ref{lemma-transversal} and its proof we see that
$I_0 = (f_1, \ldots, f_a)$ and $J_0 = (g_1, \ldots, g_b)$
where $f_1, \ldots, g_b$ is a regular system of parameters
for the regular local ring $A$. Since $I + J = \mathfrak m$, the map
$$
I \oplus J \to
\mathfrak m/\mathfrak m^2 =
\kappa f_1
\oplus \ldots \oplus
\kappa f_a
\oplus
\kappa g_1
\oplus \ldots \oplus
\kappa g_b
$$
is surjective. We conclude that we can find
$f_1', \ldots, f_a' \in I$ and $g'_1, \ldots, g_b' \in J$
whose residue classes in $\mathfrak m/\mathfrak m^2$ are equal to the
residue classes of $f_1, \ldots, f_a$ and $g_1, \ldots, g_b$.
Then $f'_1, \ldots, g'_b$ is a regular system of parameters of $A$.
By Algebra, Lemma \ref{algebra-lemma-regular-ring-CM} we find that
$A/(f'_1, \ldots, f'_a)$ is a regular local ring of dimension $b$.
Thus any nontrivial quotient of $A/(f'_1, \ldots, f'_a)$
has strictly smaller dimension
(Algebra, Lemmas \ref{algebra-lemma-regular-domain} and
\ref{algebra-lemma-one-equation}). Hence $I = (f'_1, \ldots, f'_a) = I_0$.
By symmetry $J = J_0$. This proves (2), (3), and (4).
Finally, the coefficient of $T$ in $[Y]_r \cdot [Z]_s$
is the coefficient of $T$ in $Y_{i_0} \cdot Z_{j_0}$ which is
$1$ by Lemma \ref{lemma-transversal}.
\end{proof}

\end{document}
